\documentclass{book}

%%%%%%%%%%%%%%%%%%%%%%%%%%%%%%%%%%%%%%%%%%%%%%%%%%%%%%%%%%%%%%%%%%%%%%%%%%
\usepackage{graphicx} % Required for inserting images
\usepackage{amsmath} % basic math stuff, also aligned environment 
\usepackage{amsfonts} % defeines blackboard bold fonts for \Z, \R
\usepackage{tikz} % Drawing environment 
\usepackage{amsthm} % AMS theorem defining stuff
\usepackage{amssymb} % good looking empty set
\usepackage{leftindex} % left superscript
\usepackage{float} % picture position
\usepackage{scalerel} % parallel for subscript

% Formating stuff
\usepackage{hyphenat} % forbid use of hyphen
%complete elimination of hyphen 
\righthyphenmin = 100
\lefthyphenmin = 100
%取消段前空格
\setlength{\parindent}{0pt}

% Makes section headings and cross references act like links
\usepackage[colorlinks,unicode]{hyperref}

%%%%%%%%%%%%%%%%%%%%%%%%%%%%%%%%%%%%%%%%%%%%%%%%%%%%%%%%%%%%%%%%%%%%%%%%%%
% Create simple shortcut commands.  Almost everyone who uses LaTeX
% creates commands like thes.
\newcommand{\R}{\mathbb{R}}
\newcommand{\Q}{\mathbb{Q}}
\newcommand{\Z}{\mathbb{Z}}
\newcommand{\N}{\mathbb{N}}
\newcommand{\C}{\mathbb{C}}

% this is to highlight words that are being defined and enter.
\newcommand{\define}[1]{\textbf{#1}}

% Symmetric inclusion symbol 
\newcommand{\syin}{$\subset$\kern-0.48em\raisebox{.20ex}{\tiny S}$\,$}
%\kern 左右移动(-为向左)，\raisebox 上下移动(-为向下)
% Note this symbol does not work in math mode.

% Zig-Zag product symbol 
\newcommand{\zigzag}{$\bigcirc$\kern-0.77em\raisebox{-.16ex}{\small Z}$\,$}
%\kern 左右移动(-为向左)，\raisebox 上下移动(-为向下)
% Note this symbol does not work in math mode.

% diam
\newcommand{\diam}[1]{\textrm{diam}( #1 )}
% inner product 
\newcommand{\inner}[1]{\langle #1 \rangle}
% norm 
\newcommand{\norm}[1]{\lVert #1 \rVert}
% dist
\newcommand{\dist}[1]{\textrm{dist}( #1 )}
% absolute value 
\newcommand{\abs}[1]{\lvert #1 \rvert }
% Cayley graph 
\newcommand{\Cay}[1]{\textrm{Cay}( #1 )}
% Cos graph
\newcommand{\Cos}[1]{\textrm{Cos}( #1 )}
% Automorphism group 
\newcommand{\Aut}[1]{\textrm{Aut}( #1 )}
% multiplicity of edges
\newcommand{\mul}[1]{\textrm{mul}( #1 )}
% trace of matrix
\newcommand{\tr}[1]{\textrm{tr}( #1 )}



% 简写section title


%%%%%%%%%%%%%%%%%%%%%%%%%%%%%%%%%%%%%%%%%%%%%%%%%%%%%%%%%%%%%%%%%%%%%%%%%%%%
% The next page or two of stuff is devoted to creating and
% manipulating the environments for lemma, proof, example, etc.
% The first part is pretty simple, and almost everyone has commands
% like this in every paper, article, book, etc.  
% 
% We start with things like lemmas, theorems, etc.

%注释
%\newtheorem{env name}{text}[parent counter]
%\newtheorem{env name}[shared counter]{text}
%parent counter: numbering will resteart whenever that sectional level is encountered
%shared counter: if specificed, numbering will progress sequentially for 
%all theorem element using this counter

\theoremstyle{definition} 
%adds extra space above and below, but sets the text in roman
%recommended for definitions, conditions, problems, and examples

\newtheorem{lemma}{Lemma}[chapter]
\newtheorem{proposition}[lemma]{Proposition}
\newtheorem{theorem}[lemma]{Theorem}
\newtheorem{corollary}[lemma]{Corollary}
\newtheorem{definition}[lemma]{Definition}

%自定义新的theorem environment
\newtheoremstyle{remarkstyle}
    {\topsep}%   space above, empty for `usual value'
    {\topsep}%   space below
    {}%          body font, input \itshape for italian
    {}%          indent amount 
    {}%          thm head font 
    {}%          Punctuation after thm head
    {\newline}%  Space after thm head: \newline = linebreak
    {}%          Thm head spec
\theoremstyle{remarkstyle}
\newtheorem*{remark}{Remark}%[chapter] 
\newtheorem*{example}{Example}%[chapter]
\newtheorem*{myproof}{Proof}%[chapter]
%this formating (* and %[chapter]) elminates numbering for theorem environment

%theorem之间的空间是手动调的，在最后加个\newline或者item[]

%%%%%%%%%%%%%%%%%%%%%%%%%%%%%%%%%%%%%%%%%%%%%%%%%%%%%%%%%%%%%%%%%%%%%%%%%%%%%%
%%%%%%%%%%%%%%%%%%%%%%%%%%%%%%%%%%%%%%%%%%%%%%%%%%%%%%%%%%%%%%%%%%%%%%%%%%%%%%
%%%%%%%%%%%%%%%%%%%%%%%%%%%%%%%%%%%%%%%%%%%%%%%%%%%%%%%%%%%%%%%%%%%%%%%%%%%%%%
%%%%%%%%%%%%%%%%%%%%%%%%%%%%%%%%%%%%%%%%%%%%%%%%%%%%%%%%%%%%%%%%%%%%%%%%%%%%%%
% The stuff that follows is a bit tricky, and is well more than most
% people do with LaTeX.  The basic idea is that I want to be able to
% show only examples, or only definitions, or hide only proofs, etc.  


\makeatletter
% the following key values give the set up for hiding certain
% environemnts. Note that "true" is just a dummy value.  It would work
% equally well with "foo".  Each one of these redefines the outer
% environment for an atom to turn the whole thing into a hidden
% comment.  Note that the environments theorem, example, etc. do not
% have to be written in any special way for hide and show to work.
% They get clobbered by hide.  
  \define@key{hide}{lemma}[true]{\renewenvironment{lemma}{\comment}{\endcomment}}
  \define@key{hide}{comments}[true]{\renewenvironment{comments}{\comment}{\endcomment}}
  \define@key{hide}{corollary}[true]{\renewenvironment{corollary}{\comment}{\endcomment}}
  \define@key{hide}{definition}[true]{\renewenvironment{definition}{\comment}{\endcomment}}
  \define@key{hide}{example}[true]{\renewenvironment{example}{\comment}{\endcomment}}
  \define@key{hide}{proposition}[true]{\renewenvironment{proposition}{\comment}{\endcomment}}
  \define@key{hide}{theorem}[true]{\renewenvironment{theorem}{\comment}{\endcomment}}
  \define@key{hide}{remark}[true]{\renewenvironment{remark}{\comment}{\endcomment}}
  \define@key{hide}{myproof}[true]{\renewenvironment{myproof}{\comment}{\endcomment}}
% Now this command can be given a comma separated list of names to
% hide.  Thus, \HideEnvirons{ example, proof} would hide all examples
% and proofs.
\newcommand{\HideEnvirons}[1]{\setkeys{hide}{#1}}


% The \ShowEnvirons command works like this: if it's argument is foo,
% it redefines the *hide* family key for foo to do nothing.  Then, it
% issues a \HideEnvirons command containing a list of all the atom
% types.  So, everything is hidden *unless* the hide-family key has
% been redefinied.
  \define@key{show}{lemma}[true]{\define@key{hide}{lemma}{}}
  \define@key{show}{comments}[true]{\define@key{hide}{comments}{}}
  \define@key{show}{corollary}[true]{\define@key{hide}{corollary}{}}
  \define@key{show}{definition}[true]{\define@key{hide}{definition}{}}
  \define@key{show}{example}[true]{\define@key{hide}{example}{}}
  \define@key{show}{proposition}[true]{\define@key{hide}{proposition}{}} 
  \define@key{show}{theorem}[true]{\define@key{hide}{theorem}{}}
  \define@key{show}{remark}[true]{\define@key{hide}{remark}{}}
  \define@key{show}{myproof}[true]{\define@key{hide}{myproof}{}}
% This command can be given a comma separated list of names to show,
% and only these will be shown.  Thus,
% \ShowEnvirons{example,definition} will show only examples and definitions.
\newcommand{\ShowEnvirons}[1]
{\setkeys{show}{#1}\HideEnvirons{%
    comments,
    corollary,
    definition,
    example,
    proposition,
    theorem,
    lemma,
    remark,
    myproof
  }}
\makeatother

% Comment out the next line to see the full text
\ShowEnvirons{definition, proposition, theorem, remark, example, lemma, myproof, corollary}


%%%%%%%%%%%%%%%%%%%%%%%%%%%%%%%%%%%%%%%%%%%%%%%%%%%%%%%%%%%%%%%%%%%%%%%%%%%%%%
\begin{document}

\setcounter{chapter}{7}
\chapter{Kazhdan Constant}
% Tips facing very long chapter titles 
% \chapter[medium-length title for Table of contents, if wanted]{full title name}
% \chaptermark{short title for running headers}



\section{Basics}
\begin{definition}
    Finite group $G$, $\Gamma$\syin$G$, $\rho$ a unitary representation of $G$ on representation space $V$ with $G$-invariant inner product $\inner{\cdot,\cdot}$. With nonempty $\Gamma$, define 
    $$
    \kappa (G,\Gamma,\rho,\inner{\cdot,\cdot}) = \min_{\norm{v}=1}\;\;\max_{\gamma\in\Gamma}\norm{\rho(\gamma)v-v}
    $$
    when $\Gamma$ is empty, define $\kappa$ as 0. Note that we do not require $\Gamma$ to be symmetric or a multiset.
\end{definition}
\begin{remark}
    For $G$-invariant $\inner{\cdot,\cdot}$ w.r.t. $\rho$, we have $\norm{v}=\norm{\rho (g)v}$.
\end{remark}
\begin{remark}
    Each $\norm{\rho(\gamma)v-v}$ is a continuous function of $v$, so the max of these continuous functions is a continuous function. Also, the unit $v$ lies in a compact sphere, hence minimum can be achieved. \newline
\end{remark}

\begin{proposition}
    $G,\Gamma,\rho$ as above, let $\inner{\cdot,\cdot}, \inner{\cdot,\cdot}'$ be two $G$-invariant inner products, then 
    $$
    \kappa (G,\Gamma,\rho,\inner{\cdot,\cdot}) = \kappa (G,\Gamma,\rho,\inner{\cdot,\cdot}') 
    $$
\end{proposition}
\begin{remark}
    That is, the quantity $\kappa (G,\Gamma,\rho) = \kappa (G,\Gamma,\rho,\inner{\cdot,\cdot})$ is independent from the choice of inner product. Also, note a $G$-invariant inner product always exists. 
\end{remark}
\begin{remark}
    From remark of Def. 8.1, we know the min is achieved. In fact, for any unitary irrep $\rho$, exists $v\in V,\gamma\in\Gamma$ s.t. $\norm{\rho(\gamma)\cdot v-v}=\kappa(G,\Gamma,\rho)$
\end{remark}
\begin{myproof}
    Using $\inner{\cdot,\cdot}$, we can decompose $V$ (Maschke's theorem)
    $$
    V = V_1 \bigoplus \cdots \bigoplus V_n
    $$
    We can do the same with $\inner{\cdot,\cdot}'$
    $$
    V = V'_1 \bigoplus \cdots \bigoplus V'_m
    $$
    But we know such decomposition is essentially unique, so we have $n=m$ and a G-invariant isomorphism $\phi$ s.t. $\phi(V_1)=V_1',\dots, \phi(V_n)=V'_n$. Define inner product $\inner{\cdot,\cdot}''$ on $V$ by 
    $$
    \inner{v,w}'' = \inner{\phi(v),\phi(w)}'
    $$
    Trivial to see this is indeed an inner product and is $G$-invariant. Note $V = V_1 \bigoplus \cdots \bigoplus V_n$ is also an orthogonal direct sum w.r.t. $\inner{\cdot,\cdot}''$. But by Prop. 6.37, we know invariant inner product is unique, so 
    $$
    \inner{v,w}'' = C_i\inner{v,w}
    $$
    whenever $v,w\in V_i$, so the map 
    $$
    v_1 + \cdots v_n \mapsto \frac{\phi(v_1)}{\sqrt{C_1}} + \cdots \frac{\phi(v_n)}{\sqrt{C_n}}
    $$
    defines a bijection from $\{\norm{v}=1\}$ to $\{\norm{v'}=1\}$ (mostly because $\phi$ is an isomorphism, so this is like change of basis). Then
    \begin{align*}
        \norm{\rho(\gamma)\cdot \sum^n_{i=1}v_i-\sum^n_{i=1}v_i}^2 
        &= \sum^n_{i=1}\norm{\rho(\gamma)\cdot v_i - v_i}^2 \\
        &\textrm{orthogonal, no intersection terms} \\
        &= \sum^n_{i=1}C^{-1}_i (\norm{\rho(\gamma)\cdot v_i - v_i}'')^2 \\
        &= \sum^n_{i=1}C^{-1}_i(\norm{\rho(\gamma)\cdot\phi(v_i)-\phi(v_i)}')^2 \\
        &\textrm{both steps by properties of inner products} \\
        &= (\norm{\rho(\gamma)\cdot \sum^n_{i=1}\frac{\phi(v_i)}{\sqrt{C_i}} - \sum^n_{i=1}\frac{\phi(v_i)}{\sqrt{C_i}} }')^2
    \end{align*}
    By the bijection from $v$ to $v'$, we are done. \newline 
\end{myproof}

\begin{lemma}
    If $\phi$ and $\rho$ are equivalent representations, then $\kappa(G,\Gamma,\phi)=\kappa(G,\Gamma,\rho)$
\end{lemma}
\begin{myproof}
    Let $\psi$ be a $G$-isomorphism from $V$ to $W$ (whose existence follows from two representations being equivalent). Let $\inner{\cdot,\cdot}$ be a $G$-invariant on $W$, define 
    $$
    \inner{v_1,v_2}'=\inner{\psi(v_1),\psi(v_2)}
    $$
    whose invariance follows from $G$-isomorphism $\psi$. Note $\norm{v}'=1$ iff $\norm{\psi(v)}=1$, so 
    \begin{align*}
        \kappa(G,\Gamma,\phi) 
        &= \min_{\norm{v}'=1} \max_{\gamma\in\Gamma}\norm{\phi(\gamma)v-v}' \\
        &= \min_{\norm{v}'=1} \max_{\gamma\in\Gamma}\norm{\psi(\phi(\gamma)v-v)} \\
        &= \min_{\norm{v}'=1} \max_{\gamma\in\Gamma}\norm{\rho(\gamma)\psi(v)-\psi(v) } \\
        &= \min_{\norm{w}=1}\max_{\gamma\in\Gamma}\norm{\rho(\gamma)w-w} = \kappa(G,\Gamma,\rho)
    \end{align*}
    The key is to use $\psi$ as a G-invariant isomorphism. \newline 
\end{myproof}

\begin{definition} 
    Finite nontrivial group $G$, $\Gamma$\syin$G$, define 
    $$
    \kappa(G,\Gamma) = \min_{\rho}\{\kappa(G,\Gamma,\rho)\}
    $$
    where the min is over all irreducible, nontrivial, unitary representations of $G$. This quantity is the Kazhdan constant. \newline 
\end{definition}
\begin{remark}
    Recall $G$ has only finitely many irreps up to equivalence. Also, the sum of all nontrivial irreps of $G$ is the orthogonal complement of the trivial one, so nontrivial irreps indeed exist. Hence, this min always exists.
\end{remark}
\begin{remark}
    \begin{itemize}
        \item[]
        \item One possible interpretation: given any nontrivial unitary irrep $\rho$ and any unit vector $v$ in the representation space, there exists an element in $\Gamma$ that moves $v$ a distance of at least $\kappa$, and $\kappa$ is the largest number for the preceding statement to be true. 
        \item Another implication: Let $\rho\;:\;G\rightarrow V$ be unitary irrep, $v\in V$, then exists $\gamma\in\Gamma$ s.t. 
        $$
        \norm{\rho(\gamma)v-v}\ge \kappa(G,\Gamma)\norm{v}
        $$
        \begin{itemize}
            \item Make this a unit vector by dividing $\norm{v}$
            \item Take $\gamma=\arg\max_{\gamma\in\Gamma}\norm{\rho(\gamma)w-w} $. LHS did not take min while RHS did, done.
        \end{itemize}
        \item Note 
        $$
        \kappa(G,\Gamma) = \min_{\rho}\min_{\norm{v}=1}\max_{\gamma\in\Gamma}\norm{\rho(x)v-v}
        $$
        \item[] 
    \end{itemize}
\end{remark}

\begin{proposition}
    Finite nontrivial group $G$, $\Gamma$\syin$G$, then $\kappa(G,\Gamma)\le 2$
\end{proposition}
\begin{myproof}
$$
\norm{\pi(\gamma)v-v} \le \norm{\pi(v)}+\norm{v} = 1+1 = 2
$$\newline 
\end{myproof}

\begin{proposition}
    Let $\rho\,:\,G\rightarrow GL(V)$ be a unitary representation of a finite nontrivial group $G$, let $\Gamma,\Gamma'\subset G$. Suppose that every element of $\Gamma'$ can be expressed as a word in $\Gamma$ of length less than $n$, then 
    \begin{itemize}
        \item $\kappa(G,\Gamma,\rho)\ge \kappa(G,\Gamma',\rho)/n$
        \item $\kappa(G,\Gamma)\ge \kappa(G,\Gamma')/n$ 
    \end{itemize}
\end{proposition}
\begin{myproof}
    (2) follows directly from (1). Let $v,\gamma$ be s.t. $\kappa(G,\Gamma,\rho)=\norm{\rho(\gamma)v-v} \equiv \norm{\gamma v-v}$ (for notational simplicity). Let $\gamma'\in\Gamma'$, then $\gamma'=\gamma_1\cdots\gamma_j$. Define $\alpha_k=\gamma_1\cdots\gamma_k$ for $k=1,\dots,j$, then 
    \begin{align*}
        \norm{\gamma' v-v} &= \norm{\alpha_jv-v} \\
        &= \norm{(\alpha_jv-\alpha_{j-1}v) + (\alpha_{j-1}v-\alpha_{j-2}v) + \cdots + (\alpha_2v-\alpha_{1}v) + (\alpha_1v-v)  }  \\
        &\le \norm{\alpha_{j}v-\alpha_{j-1}v} + \cdots + \norm{\alpha_{1}v-v} \\
        &= \norm{\gamma_jv-v} + \cdots \norm{\gamma _1v-v} \\
        &\textrm{G-invariant} \\
        &\le j\kappa(G,\Gamma,\rho) \\
        &\le n\kappa(G,\Gamma,\rho) 
    \end{align*}
    \newline 
\end{myproof}





\section{Kazhdan Constant, Isoperimetric Constant, and the Spectral Gap}
Denote right regular representation as $R$, ``the" trivial representation as 1. \newline 

\begin{lemma}
    Let $G$ be a finite nontrivial group, $\Gamma$\syin$G$, $d=\abs{\Gamma},\kappa=\kappa(G,\Gamma)$
    \begin{itemize}
        \item $\rho\,:\,G\rightarrow GL(V)$ be a nontrivial representation (i.e. $\inner{\chi_1,\chi_\rho}=0$, then $\kappa(G,\Gamma,\rho)\ge \frac{\kappa}{\sqrt{d}}$
        \item Let $\hat{R}$ be the restriction of $R$ to $L^2_0(G)$, $\hat{\kappa}=\kappa(G,\Gamma,\hat{R})$, then $\kappa\ge\hat{\kappa}\ge\frac{\kappa}{\sqrt{d}}$
    \end{itemize}
\end{lemma}
\begin{myproof}
    For (2), RHS is (1) on some nontrivial irrep (which must be a part of $L^2_0(G)$), done. For LHS, 
    \newline 
\end{myproof}

\begin{proposition}
    Let $G$ be a finite nontrivial group, $\Gamma$\syin$G$, $d=\abs{\Gamma},\kappa=\kappa(G,\Gamma)$, let $h=h(\Cay{G,\Gamma})$, then 
    $$
    h \ge \frac{\kappa^2}{4d}
    $$
\end{proposition}
\begin{myproof}

\end{myproof}

\begin{proposition}
    Let $G$ be a finite nontrivial group, $\Gamma$\syin$G$, $d=\abs{\Gamma},\kappa=\kappa(G,\Gamma)$, let $\lambda_1$ be the second largest eigenvalue, then 
    $$
    \kappa \ge \sqrt{\frac{2(d-\lambda_1)}{d}}
    $$
\end{proposition}
\begin{myproof}
    
\end{myproof}

\begin{theorem}
    Let $d$ be a positive integer, let $\{G_n\}$ be a sequence of groups with $\abs{G_n}\rightarrow\infty$. For each $n$, let $\Gamma_n$\syin $G_n$ s.t. $\abs{\Gamma_n}=d$, then $\{\Cay{G_n,\Gamma_n}\}$ is an expander family iff there exists $\epsilon>0$ s.t. $\kappa(G_n,\Gamma_n)\ge\epsilon$ for all $n$.
\end{theorem}
\begin{myproof}
    With Prop.8.8 and Prop.8.9, we get a representation copy of Chapter 1. \newline 
\end{myproof}

\begin{corollary}
    Let $G$ be a finite nontrivial group, then $\kappa(G,G)>\sqrt{2}$
\end{corollary}
\begin{myproof}

\end{myproof}





\section{Abelian Groups Never Yields Expander Families}
\begin{lemma}
    Let $\theta$ be a real number, then $\abs{e^{i\theta}-1}\le\abs{\theta}$
\end{lemma}
\begin{myproof}
    By elementary calculus. \newline 
\end{myproof}

\begin{proposition}
    Let $G$ be a nontrivial finite abelian group, $\Gamma\subset G$, $d=\abs{\Gamma}$, then 
    $$
    \kappa(G,\Gamma)\le \frac{2\pi}{\abs{G}^{1/d}-1}
    $$
\end{proposition}
\begin{remark}
    We see $\kappa \rightarrow 0$ as $\abs{G}\rightarrow\infty $, so sequence abelian groups cannot yield expander family.
\end{remark}
\begin{myproof}
    First, recall the irreps of finite abelian groups \newline 
    Suppose $G=\Z_{n_1}\times \cdots \Z_{n_r} $, then the set of irreps of $G$ is
    $$
    \{\phi_a\,\lvert\, a = (a_1,\dots,a_r) \in \Z_{n_1}\times \cdots \Z_{n_r}\}
    $$
    where 
    $$
    \phi_a(k_1,\dots,k_r) = (e^{\frac{2\pi ia_1k_1}{n_1}}\dots e^{\frac{2\pi ia_rk_r}{n_r}})
    $$
    Note $\abs{G}=\abs{\Z_{n_1}}\times\cdots\abs{\Z_{n_r}}$, so the size of set of all irreps is the same as $\abs{G}$ \newline 

Let $\kappa = \kappa(G,\Gamma) >0$, $C$ be the smallest integer greater than or equal to $2\pi /\kappa$. Claim $\abs{G}\le C^d$. If we have this, then $\abs{G}^{1/d} \le C\le 2\pi/\kappa +1$, solve for $\kappa$, done. \newline 

Suppose $\abs{G}>C^d$. Let $G=\Z_{n_1}\times \cdots \Z_{n_r}$, $S$ be the set of all irreps. We know that for any irrep $\rho$, for arbitrary $g\in G$, $\rho(g)$ is of the form $e^{i\theta} $ for some $\theta \in [0,2\pi)$. Note the image of $\rho$ will always lie inside $\C^*$ ($\C$ without 0). \newline 

For any element $z$ on the unit circle in $\C$, let $\ell(z)$ be the unique argument angle s.t. $e^{i\theat} = z $. Define $\phi\,:\, S\rightarrow [0,2\pi)^d$ by 
$$
\phi(\rho) = [\ell(\rho(\gamma_1)), \dots, \ell(\rho(\gamma_n))]
$$
where $\gamma_1,\dots,\gamma_d$ are the elements of $\Gamma$. That is, we take the argument for the mapped results of $\Gamma$ \newline 

Suppose $a_1,\dots,a_d$ are integers with $1\le a_j\le C$ for all $j$, and $a=(a_1,\dots, a_d)$, define 
$$
B_a = \left\{ 
(r_1, \dots, r_d) \in [0,2\pi)^d \;:\; 
\frac{2\pi(a_j-1)}{C} \le r_j \le \frac{2\pi a_j}{C}
\right\}
$$
That is, we partition the $[0,2\pi)^d$. Note there are $C^d$ choices of $B_a$, and the union of them all indeed partitions $[0,2\pi)^d$. The figure below shows the case for $C=3, d=2$, the x-axis denotes $r_1$, the y-axis denotes $r_2$, and $B_{(1,2)} $ is the case $a_1 = 1, a_2 = 2$. \newline 

We know $\abs{S}=\abs{G}$, by the pigeonhole principle, there must be two distinct irreps $\rho_1,\rho_2$ falling in the same $B_a$ box. That is, we can control their corresponding arguments
$$
\omega \equiv \abs{\ell(\rho_1(\gamma_j)) - \ell(\rho_2(\gamma_j)) } < \frac{2\pi}{C}
$$
for any $j$. \newline 

Define $\rho\,:\, G\rightarrow GL(1,\C)$ by $\rho(g) = \rho_1(g)/\rho_2(g)$. Straightforward to see this is a homomorphism (hence a representation). Also, $dim(GL(1,\C))=1$ and $\rho_1(g)/\rho_2(g)\neq 1$, so it is a nontrivial irrep. Note $\rho(\gamma_j) = e^{i\omega} $, so by Lemma 8.12, $\abs{\rho(\gamma_j)-1}<2\pi/C$ for all $j$. Note the unit vector in $\C$ is exactly 1, so $\kappa < \kappa(G,\Gamma,\rho) < 2\pi\C$, contradicts with definition of $C$, done. \newline 

\end{myproof}





\section{Kazhdan Constant, Subgroups, and Quotients}
\begin{definition}
    Groups $G,H$, $\rho\,:\, H\rightarrow GL(V)$ a representation, $\phi\,:\, G\rightarrow H$ a homomorphism. Let $\rho* = \rho \,  \circ \, \phi $, then $\rho^*\,:\, G\rightarrow GL(V)$ is a representation of $G$, called the pullback of $\rho$ via $\phi$ \newline 
\end{definition}

\begin{lemma}
    Finite group $G$, $N\vartriangleleft G$, suppose $\rho\,:\, G\rightarrow GL(V)$ is an irrep of $G/N$, let $\phi\,:\, G\rightarrow G/N$ be the canonical homomorphism (i.e. $\phi(g) = gN$), let $\rho^*$ be the pullback of $\rho$ via $\phi$, then $\rho^*$ is an irrep of $G$.
\end{lemma}
\begin{myproof}
    Homomorphism is trivial. Suppose there is a proper nontrivial subspace $W$ of $V$ s.t. $\rho^*(g)w \in W$, then $\rho(\phi(g))w\in W$, thus $\rho(gN)w\in W$, contradicting with irreducibility of $\rho$. Essentially, any subspace invariant under action of $G$ is invariant under the action of $G/N$.
\end{myproof}
\begin{remark}
    If an inner product is $G/N$-invariant under $\rho$, then it is also $G$-invariant under the pullback $\rho^*$. \newline 
\end{remark}

\begin{proposition}
    Finite nontrivial group $G$, finite nontrivial $N\vartriangleleft G$, $\Gamma$\syin $G$. Let $\bar{\Gamma}$ be the image of $\Gamma$ under the canonical homomorphism, then 
    $$
    \kappa(G,\Gamma) \le \kappa(G/N, \bar{\Gamma})
    $$
\end{proposition}
\begin{myproof}
    Let $\rho\,:\,G/N\rightarrow GL(V)$ be a nontrivial irrep. Let unit vector $v\in V$, $\gamma_0N\in \bar{\Gamma}$ be s.t. $\norm{\rho(\gamma_0N)v-v}=\kappa(G/N,\bar{\Gamma},\rho)$ (we can do this since the min is achievable). Let $\rho^*$ be the pullback of $\rho$ via the canonical homomorphism, then for any $\gamma\in\Gamma$, 
    $$
    \norm{\rho^*(\gamma)v-v} = \norm{\rho(\phi(\gamma))v-v} = \norm{\rho(\gamma N)v-v}
    $$
    So 
    $$
    \kappa(G,\Gamma) \le \kappa(G,\Gamma,\rho^*) = \kappa(G/N,\bar{\Gamma},\rho)
    $$
    this works for any $\rho$ of $G/N$, done. 
\end{myproof}
\begin{remark}
    With Prop. 8.16 and Theorem 8.10, we recover the Quotients Nonexpansion Principle. \newline 
\end{remark}

\begin{proposition}
    Let $G$ be a finite nontrivial subgroup, $H$ be a nontrivial subgroup of $G$, $\Gamma$\syin $G$, $d=\abs{\Gamma}, n=[G:H], T=\{t_1,\dots,t_n\}$, where $T$ is the set of transverals for $H$ in $G$. Let $\hat{\Gamma}$ be the set of Schreier generators. Then 
    $$
    \kappa(G,\Gamma) \le d^{1/2}\kappa(H,\hat{\Gamma})
    $$
\end{proposition}
\begin{myproof}
    Recall some useful facts:
    \begin{itemize}
        \item $g\in G$, $g = h\overline{g}$
        \item $\widehat{(t,\gamma)} = t\gamma(\overline{t\gamma})^{-1} $
        \item Any $g\in G$, $g=ht$ for some unique $h\in H, t\in T $, and $t\gamma = \widehat{(t,\gamma)}\overline{t\gamma} $
    \end{itemize}
\end{myproof}

Let $\sigma\,:\,H\rightarrow GL(W)$ be a nontrivial irrep of $H$, $\inner{\cdot,\cdot}$ be an $H$-invariant innner product on $W$, $w$ be the unit vector in $W$ s.t. 
$$
\kappa(H,\hat{\Gamma}) = \max_{\widehat{(t,\gamma)}\in\hat{\Gamma}} \norm{\sigma(\widehat{(t,\gamma)})w -w}_W
$$
Let $\rho=Ind^G_H(\sigma)$ be the induced representation, $V=\{f\,:\, G\rightarrow W \,\lvert\, f(hg) = \sigma(h)f(g), \forall g\in G, h\in H\} $. \newline 

For all $h\in H, t\in T$, define 
$$
f(ht) = n^{-1/2}\sigma(h)w
$$
Note $f(h'g) = f(h'h''t) = f(ht) = n^{-1/2}\sigma(h)w  $ and $\sigma(h')f(g) = \sigma(h')f(h''t) = \sigma(h'h'')n^{-1/2}w $, so indeed $f\in V$. Also, $f(t) = n^{-1/2}w$. \newline 

Follow the inner product 
$$
\inner{y_1, y_2}_V = \sum^{n}_j\inner{y_1(t_1), y_2(t_1)}_W
$$
we have 
$$
\norm{f}^2_V = \sum^n_{j=1}\inner{n^{-1/2}w, n^{-1/2}w}_W = \sum^n_{j=1}n^{-1}\norm{w}^2_W = 1
$$
so $f$ is a unit vector in $V$, which can be used to get bounds for $\kappa$ \newline 

Let $\gamma\in \Gamma$, 
\begin{align*}
    \norm{\rho(\gamma)f - f}^2_V 
    &= \sum^n_{j=1}\norm{f(t_j\gamma)-f(t_j)}^2_W \\
    &\textrm{the definition of Ind, $\rho(x)f(g) = f(gx)$} \\
    &= \sum^n_{j=1} \norm{f(\widehat{(t_j,\gamma)}\overline{t_j\gamma}) - f(t_j) }^2_W \\
    &= \sum^n_{j=1} \norm{\sigma\left(\widehat{(t_j,\gamma)}\right) f(\overline{t_j\gamma}) - f(t_j) }^2_W \\
    &= \sum^n_{j=1} \norm{\sigma\left(\widehat{(t_j,\gamma)}\right) (n^{-1/2}w) - n^{-1/2}w }^2_W \\
    &= \frac{1}{n} \sum^n_{j=1} \norm{\sigma\left(\widehat{(t_j,\gamma)}\right) w - w }^2_W
\end{align*}
Then for some $j$ (i.e. the max of all these terms)
$$
\norm{\rho(\gamma)f - f}^2_V \le \norm{\sigma\left(\widehat{(t_j,\gamma)}\right) w - w }^2_W
$$
So 
$$
\norm{\rho(\gamma)f - f}_V \le \kappa(H,\hat{\Gamma})
$$
Since $\gamma$ is arbitrary, we have $\kappa(G,\Gamma,\rho)\le \kappa(H,\hat{\Gamma})$. Also, we know from the last example of Chapter 6, $\rho$ cannot contain the trivial representation, so by Lemma 8.7, $\kappa(G,\Gamma)/d^{1/2} \le \kappa(G,\Gamma,\rho) \le \kappa(H,\hat{\Gamma})  $, done. 
\end{myproof}
\begin{remark}
    With Prop. 8.17 and Theorem 8.10, we recover the Subgroups Nonexpansion Principle.
\end{remark}













\end{document}
